\section*{Chapitre \ref{ch:g4:what-plants-need} --- Ce dont les plantes ont besoin}

\begin{solution}{exo:g4:what-plants-need:1}
L'eau, la lumière, l'air, la chaleur, et les \omterm{prop:g4:what-plants-need:needs}{sels minéraux} (bus
dissous dans l'eau du sol).
\end{solution}

\begin{solution}{exo:g4:what-plants-need:2}
Exactement une --- la chose testée~; tout le reste doit être
identique. Le pot auquel on ne touche pas est le témoin.
\end{solution}

\begin{solution}{exo:g4:what-plants-need:3}
Jaune pâle au lieu de vert, étrangement grand et étiré, faible ---
puis elle renonce. La lumière manquait, si bien que la \omterm{def:g1:plants-around-us:plant}{plante} ne
pouvait fabriquer ni sa substance verte, ni sa nourriture.
\end{solution}

\begin{solution}{exo:g4:what-plants-need:4}
De toutes petites quantités de substances nourrissantes présentes
dans le sol. Les \omterm{def:g1:plants-around-us:parts}{racines} les boivent dissoutes dans l'eau du sol.
\end{solution}

\begin{solution}{exo:g4:what-plants-need:5}
Parce que les récoltes emportent les \omterm{prop:g4:what-plants-need:needs}{sels minéraux} du sol année
après année~; le compost les réapprovisionne. Il nourrit le
\emph{sol}, et l'eau du sol porte ensuite les sels jusqu'aux
\omterm{def:g1:plants-around-us:parts}{racines}.
\end{solution}

\begin{solution}{exo:g4:what-plants-need:6}
Un \omterm{def:g1:animals-around-us:animal}{animal} grandit à partir des \omterm{def:g1:living-or-not:living}{êtres vivants} qu'il mange~; une
\omterm{def:g1:plants-around-us:plant}{plante} ne mange personne --- elle fabrique sa propre matière à
partir d'eau, de \omterm{prop:g4:what-plants-need:needs}{sels minéraux}, d'air et de l'\omterm{prop:g3:balanced-diet:jobs}{énergie} de la lumière
captée par ses \omterm{def:g1:plants-around-us:parts}{feuilles} vertes.
\end{solution}

\begin{solution}{exo:g4:what-plants-need:7}
Pas du sol~: de l'eau, de l'air, la lumière fournissant l'\omterm{prop:g3:balanced-diet:jobs}{énergie}
--- le tout bâti dans les \omterm{def:g1:plants-around-us:parts}{feuilles} vertes. (Le sol n'a apporté que
de l'eau et de petites quantités de \omterm{prop:g4:what-plants-need:needs}{sels minéraux}.)
\end{solution}

\begin{solution}{exo:g4:what-plants-need:8}
Réponse libre. Attendu~: le pot arrosé vert et droit, son jumeau sec
fané en quelques jours --- et comme seule l'eau différait, c'est
l'eau qui en est responsable.
\end{solution}

\begin{solution}{exo:g4:what-plants-need:9}
Il n'y a pas de témoin~: un seul pot, et tout à la fois --- ce pot
avait peut-être simplement une bonne lumière et un arrosage assidu
(les chanteurs rendent souvent visite à leurs \omterm{def:g1:plants-around-us:plant}{plantes}~!). La version
honnête~: deux pots identiques, mêmes \omterm{prop:g3:flowers-fruits-seeds:transform}{graines}, même terre, même
fenêtre, même arrosage~; chanter à un seul~; comparer au bout d'un
mois. Alors seulement une différence désignerait le chant.
\end{solution}

\begin{solution}{exo:g4:what-plants-need:10}
La lumière~: elle pousse immergée, mais dans une eau éclairée, ses
\omterm{def:g1:plants-around-us:parts}{feuilles} vertes captant encore la lumière à travers elle. L'eau qui
l'entoure fournit aussi ce que fournit l'eau du sol --- des \omterm{prop:g4:what-plants-need:needs}{sels
minéraux} dissous, bus par toute sa surface au lieu de l'être par des
\omterm{def:g1:plants-around-us:parts}{racines}.
\end{solution}

\begin{solution}{exo:g4:what-plants-need:11}
Les \omterm{prop:g4:what-plants-need:needs}{sels minéraux}~: l'eau, la lumière, l'air et la chaleur étaient
égaux, et du sable pur arrosé d'eau de pluie n'apporte presque aucun
sel. Elle a poussé au début sur les réserves rangées dans ses
propres tissus (comme la \omterm{prop:g2:from-seed-to-plant:cycle}{jeune pousse} enterrée poussait sur la
réserve de la \omterm{def:g2:from-seed-to-plant:seed}{graine}) --- les réserves dépensées, le manque s'est
montré.
\end{solution}
